% !TEX root = ../SANER2027-main-labeling.tex
% Conclusion
% Claims to support: summary + future work

\section{Conclusion and Future Work}
\label{sec:conclusion}

\begin{comment}
We evaluate all three against a LoRA fine-tuning baseline on an 11-project benchmark, in both \emph{project-specific} (PS) and \emph{project-agnostic} (PA) data scope settings. Experiments use Qwen2.5-Instruct at four sizes.
In this study, we investigated retrieval-augmented few-shot prompting as a cost-efficient alternative to LoRA fine-tuning for issue report classification (IRC). Given a query issue, we retrieved its top-$k$ nearest neighbors from historical issue reports to be used in three proposed methods. First, \rag, which presents the neighbors as labeled examples in a few-shot prompt. Second, \frag, a \rag\ variant that drops bug-labeled neighbors when the retrieved label distribution indicates the query is likely a question. Third, \knn, which predicts the label by a weighted vote on the neighbors' labels, establishing a non-parametric baseline. 

We evaluated all three methods against the LoRA fine-tuning baseline using the Qwen2.5-Instruct model at four sizes. Evaluations were conducted under both project-specific (PS) and project-agnostic (PA) data scope settings.
The results show that retrieval-augmented few-shot prompting is a practical alternative to fine-tuning for IRC. \rag\ PS outperforms fine-tuning PS by 0.021--0.040 macro $F_1$ across all four model sizes, and trails fine-tuning PA by 0.029 macro $F_1$ in aggregate despite using 11$\times$ less labeled data per project. \frag\ PS closes this gap to TOST equivalence at $\delta{=}0.02$, tightening to statistical equivalence at $\delta{=}0.01$ under the \knn\ fallback protocol for invalid LLM outputs. Averaged across all models, retrieval-augmented few-shot also required 33\% less peak GPU memory than fine-tuning.

These findings position retrieval-augmented few-shot prompting as the more practical IRC method, particularly when hardware is constrained, labeled data is limited, or deployment requires frequent model or data updates. We encourage future work to evaluate these methods in diverse data distribution scenarios and to explore additional retrieval techniques and configurations.

%%%%%%% Future Work
Future work should examine the generalizability of our findings beyond the studied benchmark by evaluating IRC methods under more diverse conditions, including imbalanced label distributions, minority-class scenarios, multi-label classification, and other issue tracking platforms such as Jira and Bugzilla. Additional work should also explore stronger retrieval-side interventions, such as LLM-based filtering or re-ranking of retrieved neighbors, as well as the effects of different embedding models and similarity metrics on retrieval quality. Finally, future studies should investigate alternative prompt structures, including advanced strategies such as Chain-of-Thought prompting~\cite{koyuncu2025exploring}, and evaluate other open-source LLM families to better understand how model architecture influences retrieval-augmented few-shot IRC.
\end{comment}

In this study, we investigated retrieval-augmented few-shot prompting (\rag) as a training-free alternative to LoRA fine-tuning for issue report classification (\irc). We conducted an empirical study with four Qwen2.5-Instruct models and 6{,}600 issue reports from eleven projects, where the retrieval index and the fine-tuning data come either from the target project's 300 labeled issues or from all 3{,}300 labeled issues of the eleven projects. We first tested how well the labels of a query issue's most similar neighbors predict its label, using a similarity-weighted vote over them (\knn). Since these neighbors carry label information, \rag gives them to the LLM as few-shot examples. Because \rag\ most often labels questions as bugs, a filtered variant (\frag) removes the bug-labeled examples for suspected questions.

Retrieved examples improve the LLM over both zero-shot prompting and \knn, and \frag\ improves on \rag\ by reducing question-to-bug errors, at the cost of bug recall. With the same labeled issues of the target project, \rag\ and \frag\ at their best number of examples outperform LoRA fine-tuning at every model size. When fine-tuning instead uses all 3{,}300 labeled issues of the eleven projects, it outperforms \rag, while its lead over \frag\ is not statistically significant. With \knn\ replacing invalid LLM outputs, \frag\ matches fine-tuning (+0.001 macro~$F_1$ on average). Fine-tuning finds more of the actual bugs but over-predicts the bug label, whereas \frag\ makes fewer false bug predictions at most model sizes, and the question-to-bug confusion persists under every method we test. Needing only the target project's labeled issues and 33\% less peak GPU memory on average, retrieval-augmented few-shot prompting is a competitive, training-free alternative to LoRA fine-tuning.

The two approaches trade different strengths, and which one suits a project depends on its labeled data, its GPU budget, and whether a missed bug or a false bug label costs it more. With retrieval, a newly labeled issue only needs to be indexed, and the LLM can be replaced without retraining. As this work constitutes a case study on the Qwen2.5-Instruct family, future work can focus on the generalizability of these findings by extending them to other LLM families and quantization settings to further strengthen or bound the conclusions drawn here. Additional directions include examining imbalanced label distributions, minority-class scenarios, multi-label classification, and other issue tracking platforms such as Jira and Bugzilla. Further research should also explore stronger retrieval-side interventions, such as LLM-based filtering or re-ranking of retrieved neighbors, as well as the effects of different embedding models and similarity metrics on retrieval quality. Finally, future studies should investigate alternative prompt structures, including advanced strategies such as Chain-of-Thought prompting~\cite{koyuncu2025exploring}.
